\begin{center}
    \textbf{\LARGE ĐỀ KIỂM TRA ĐÁNH GIÁ GIỮA KÌ I}\\
    \textbf{\Large MÔN: HÓA HỌC 10 -- CHƯƠNG TRÌNH GDPT 2018}\\
    \textbf{\large BỘ SÁCH CÁNH DIỀU / KẾT NỐI TRI THỨC / CHÂN TRỜI SÁNG TẠO}\\
    \textit{Thời gian làm bài: 45 phút, không kể thời gian phát đề}
\end{center}
\hrule
\vspace{0.5cm}

\noindent\textbf{\large PHẦN I. Câu trắc nghiệm nhiều phương án lựa chọn.} Thí sinh trả lời từ câu 1 đến câu 16. Mỗi câu hỏi thí sinh chỉ chọn một phương án.

\begin{ex}
    Hóa học là ngành khoa học thuộc lĩnh vực khoa học tự nhiên. Đối tượng nghiên cứu chính của hóa học là
    \choice
    {sự chuyển động cơ học của các thiên thể trong hệ Mặt Trời}
    {\True thành phần, cấu trúc, tính chất và sự biến đổi của chất cũng như ứng dụng của chúng}
    {sự sinh trưởng, phát triển và tập tính sinh học của các loài sinh vật}
    {các hiện tượng địa chấn, khí tượng và sự biến đổi của vỏ Trái Đất}
    \loigiai{
        Theo định nghĩa, hóa học là ngành khoa học thuộc lĩnh vực khoa học tự nhiên, nghiên cứu về thành phần, cấu trúc, tính chất và sự biến đổi của các chất cũng như ứng dụng của chúng trong đời sống và sản xuất.
    }
\end{ex}

\begin{ex}
    Hóa học đóng vai trò thiết yếu đối với đời sống, sản xuất và phát triển kinh tế xã hội. Ứng dụng nào sau đây thể hiện vai trò trực tiếp của ngành hóa học trong việc bảo vệ và chăm sóc sức khỏe con người?
    \choice
    {Sản xuất xi măng, gạch ngói và vật liệu xây dựng}
    {Sản xuất phân bón hóa học và thuốc bảo vệ thực vật}
    {\True Nghiên cứu và sản xuất các loại thuốc chữa bệnh, vắc-xin và dịch truyền}
    {Chế biến dầu mỏ để cung cấp nhiên liệu cho các phương tiện giao thông}
    \loigiai{
        Ngành công nghiệp dược phẩm hóa học chuyên nghiên cứu, tách chiết và tổng hợp các loại thuốc kháng sinh, thuốc điều trị, vắc-xin, vitamin và khoáng chất nhằm phòng ngừa và chữa trị bệnh tật cho con người.
    }
\end{ex}

\begin{ex}
    Để nghiên cứu sự ăn mòn kim loại sắt trong các môi trường khác nhau, một nhóm học sinh đã thực hiện các bước sau:
    \begin{itemize}
        \item[(1)] Đưa ra dự đoán: ``Đinh sắt bị gỉ nhanh nhất trong môi trường nước muối''.
        \item[(2)] Xác định câu hỏi nghiên cứu: ``Môi trường nào làm đinh sắt bị ăn mòn nhanh nhất?''.
        \item[(3)] Tiến hành thí nghiệm ngâm các đinh sắt vào các ống nghiệm chứa: nước cất, nước muối và không khí khô.
        \item[(4)] Rút ra kết luận về các yếu tố thúc đẩy quá trình ăn mòn kim loại sắt.
        \item[(5)] Quan sát, ghi chép hiện tượng gỉ sét trên bề mặt đinh sắt sau một tuần.
    \end{itemize}
    Trình tự thực hiện hợp lí nhất theo phương pháp nghiên cứu hóa học là
    \choice
    {\True $(2) \rightarrow (1) \rightarrow (3) \rightarrow (5) \rightarrow (4)$}
    {$(2) \rightarrow (3) \rightarrow (1) \rightarrow (5) \rightarrow (4)$}
    {$(1) \rightarrow (2) \rightarrow (3) \rightarrow (5) \rightarrow (4)$}
    {$(3) \rightarrow (5) \rightarrow (2) \rightarrow (1) \rightarrow (4)$}
    \loigiai{
        Quy trình nghiên cứu hóa học theo phương pháp thực nghiệm chuẩn gồm các bước:
        \begin{itemize}
            \item Bước 1: Xác định vấn đề nghiên cứu (2).
            \item Bước 2: Đề xuất giả thuyết khoa học / dự đoán (1).
            \item Bước 3: Lập kế hoạch và tiến hành thí nghiệm (3).
            \item Bước 4: Thu thập số liệu, quan sát hiện tượng (5).
            \item Bước 5: Phân tích kết quả và rút ra kết luận (4).
        \end{itemize}
        Vậy thứ tự đúng là $(2) \rightarrow (1) \rightarrow (3) \rightarrow (5) \rightarrow (4)$.
    }
\end{ex}

\begin{ex}
    Nguyên tử có cấu tạo gồm hạt nhân nằm ở tâm mang điện tích dương và lớp vỏ mang điện tích âm. Các loại hạt cấu tạo nên hạt nhân nguyên tử là
    \choice
    {proton và electron}
    {neutron và electron}
    {\True proton và neutron}
    {proton, neutron và electron}
    \loigiai{
        Hạt nhân nguyên tử được cấu tạo từ các hạt proton (mang điện tích dương) và neutron (không mang điện), ngoại trừ hạt nhân của nguyên tử $^1_1\text{H}$ chỉ gồm 1 proton. Lớp vỏ nguyên tử được tạo nên từ các electron mang điện tích âm.
    }
\end{ex}

\begin{ex}
    Trong các loại hạt cấu tạo nên nguyên tử, loại hạt mang điện tích âm và có khối lượng không đáng kể (khoảng $0{,}00055\text{ amu}$) là
    \choice
    {proton}
    {\True electron}
    {neutron}
    {hạt alpha}
    \loigiai{
        Hạt electron mang điện tích âm ($-1$ đơn vị điện tích nguyên tố, tương đương $-1{,}602 \times 10^{-19}\text{ C}$) và có khối lượng $m_e \approx 0{,}00055\text{ amu} \approx 9{,}109 \times 10^{-28}\text{ g}$, rất nhỏ so với khối lượng của proton và neutron ($m_p \approx m_n \approx 1\text{ amu}$).
    }
\end{ex}

\begin{ex}
    Nguyên tử của nguyên tố aluminium (nhôm) có $13$ hạt proton và $14$ hạt neutron trong hạt nhân. Kí hiệu nguyên tử của aluminium là
    \choice
    {$^{14}_{13}\text{Al}$}
    {\True $^{27}_{13}\text{Al}$}
    {$^{13}_{27}\text{Al}$}
    {$^{27}_{14}\text{Al}$}
    \loigiai{
        Số hiệu nguyên tử $Z = \text{số proton} = 13$.\\
        Số khối $A = Z + N = 13 + 14 = 27$.\\
        Kí hiệu nguyên tử có dạng $^A_Z\text{X}$, do đó kí hiệu của aluminium là $^{27}_{13}\text{Al}$.
    }
\end{ex}

\begin{ex}
    Các đồng vị của cùng một nguyên tố hóa học là những nguyên tử có
    \choice
    {cùng số neutron nhưng khác nhau về số proton}
    {\True cùng số proton nhưng khác nhau về số neutron}
    {cùng số khối nhưng khác nhau về số proton}
    {cùng số electron nhưng khác nhau về số proton}
    \loigiai{
        Đồng vị là những nguyên tử có cùng số đơn vị điện tích hạt nhân (cùng số proton hay số hiệu nguyên tử $Z$) nhưng có số neutron khác nhau, do đó số khối $A$ của chúng khác nhau.
    }
\end{ex}

\begin{ex}
    Khẳng định nào sau đây là \textbf{đúng} về hình dạng orbital nguyên tử (AO) và số electron tối đa trong một AO?
    \choice
    {AO s có dạng hình số tám nổi và chứa tối đa 2 electron}
    {AO p có dạng hình cầu và chứa tối đa 6 electron}
    {\True AO s có dạng hình cầu và mỗi AO chỉ chứa tối đa 2 electron}
    {AO p có dạng hình số tám nổi và mỗi AO p chứa tối đa 6 electron}
    \loigiai{
        - AO s có dạng hình cầu đối xứng tâm trong không gian; AO p có dạng hình số tám nổi (gồm $p_x, p_y, p_z$).\\
        - Theo nguyên lí Pauli, mỗi một orbital nguyên tử (dù là s, p, d hay f) chỉ chứa tối đa 2 electron và 2 electron này có chiều tự quay ngược nhau (spin trái dấu). Phân lớp p gồm 3 AO p nên chứa tối đa $3 \times 2 = 6$ electron, nhưng mỗi AO p đơn lẻ chỉ chứa tối đa 2 electron.
    }
\end{ex}

\begin{ex}
    Các electron thuộc cùng một phân lớp electron trong nguyên tử có đặc điểm nào sau đây?
    \choice
    {Có mức năng lượng gần bằng nhau}
    {\True Có mức năng lượng bằng nhau}
    {Có cùng hướng định hướng không gian của orbital}
    {Luôn có cùng chiều tự quay (spin)}
    \loigiai{
        - Các electron thuộc cùng một lớp electron có mức năng lượng \textit{gần bằng nhau}.\\
        - Các electron thuộc cùng một phân lớp electron có mức năng lượng \textit{bằng nhau}.
    }
\end{ex}

\begin{ex}
    Trong bảng tuần hoàn các nguyên tố hóa học hiện đại, số thứ tự của chu kì bằng
    \choice
    {số electron hóa trị của nguyên tử}
    {số electron ở lớp ngoài cùng của nguyên tử}
    {\True số lớp electron của nguyên tử}
    {số proton trong hạt nhân nguyên tử}
    \loigiai{
        Chu kì là dãy các nguyên tố mà nguyên tử của chúng có cùng số lớp electron, được sắp xếp theo chiều tăng dần của điện tích hạt nhân. Do đó, số thứ tự của chu kì luôn bằng số lớp electron trong nguyên tử.
    }
\end{ex}

\begin{ex}
    Trong một chu kì, đi từ trái sang phải theo chiều tăng dần của điện tích hạt nhân, tính chất acid -- base của các oxide cao nhất và hydroxide tương ứng biến đổi theo quy luật nào sau đây?
    \choice
    {\True Tính acid tăng dần, tính base giảm dần}
    {Tính acid giảm dần, tính base tăng dần}
    {Cả tính acid và tính base đều tăng dần}
    {Cả tính acid và tính base đều giảm dần}
    \loigiai{
        Trong một chu kì, từ trái sang phải theo chiều tăng điện tích hạt nhân, tính kim loại giảm dần và tính phi kim tăng dần. Tương ứng với quy luật đó, tính base của các oxide cao nhất và hydroxide giảm dần, đồng thời tính acid của chúng tăng dần.
    }
\end{ex}

\begin{ex}
    Nguyên tố phosphorus ($Z = 15$) thuộc nhóm VA trong bảng tuần hoàn. Công thức oxide cao nhất của phosphorus là
    \choice
    {$\text{P}_2\text{O}_3$}
    {$\text{PO}_2$}
    {\True $\text{P}_2\text{O}_5$}
    {$\text{PO}_3$}
    \loigiai{
        Với các nguyên tố nhóm A, hóa trị cao nhất trong oxide với oxygen bằng số thứ tự của nhóm. Nguyên tố phosphorus thuộc nhóm VA nên có hóa trị cao nhất là V trong oxide, công thức oxide cao nhất là $\text{P}_2\text{O}_5$.
    }
\end{ex}

\begin{ex}
    Hydroxide tương ứng với oxide cao nhất của nguyên tố sulfur ($Z = 16$, nhóm VIA trong bảng tuần hoàn) là
    \choice
    {$\text{H}_2\text{SO}_3$}
    {\True $\text{H}_2\text{SO}_4$}
    {$\text{S(OH)}_4$}
    {$\text{H}_2\text{S}$}
    \loigiai{
        Nguyên tố sulfur thuộc nhóm VIA nên hóa trị cao nhất trong oxide là VI (oxide cao nhất là $\text{SO}_3$). Hydroxide tương ứng của $\text{SO}_3$ là sulfuric acid có công thức $\text{H}_2\text{SO}_4$ (trong đó S có hóa trị VI).
    }
\end{ex}

\begin{ex}
    Nội dung nào sau đây phát biểu đúng về định luật tuần hoàn các nguyên tố hóa học?
    \choice
    {\True Tính chất của các nguyên tố và đơn chất, cũng như thành phần và tính chất của các hợp chất tạo nên từ chúng biến đổi tuần hoàn theo chiều tăng của điện tích hạt nhân nguyên tử}
    {Tính chất của các nguyên tố và đơn chất biến đổi tuần hoàn theo chiều tăng của nguyên tử khối tương đối}
    {Bán kính nguyên tử và khối lượng nguyên tử biến đổi tuần hoàn theo chiều tăng của số khối}
    {Số lớp electron trong nguyên tử biến đổi tuần hoàn theo chiều tăng của số lượng neutron trong hạt nhân}
    \loigiai{
        Định luật tuần hoàn (hiện đại): Tính chất của các nguyên tố và đơn chất, cũng như thành phần và tính chất của các hợp chất tạo nên từ chúng biến đổi tuần hoàn theo chiều tăng của điện tích hạt nhân nguyên tử.
    }
\end{ex}

\begin{ex}
    Nguyên nhân sâu xa dẫn tới sự biến đổi tuần hoàn tính chất của các nguyên tố hóa học khi điện tích hạt nhân tăng dần là do sự biến đổi tuần hoàn của
    \choice
    {số neutron trong hạt nhân nguyên tử}
    {khối lượng nguyên tử tương đối}
    {\True cấu hình electron lớp ngoài cùng của nguyên tử}
    {tỉ lệ giữa số hạt mang điện và không mang điện trong nguyên tử}
    \loigiai{
        Cấu hình electron lớp ngoài cùng quyết định tính chất hóa học cơ bản của nguyên tố. Sự lặp lại tuần hoàn cấu hình electron lớp ngoài cùng của nguyên tử các nguyên tố theo chiều tăng dần của điện tích hạt nhân chính là nguyên nhân dẫn đến sự biến đổi tuần hoàn tính chất của các đơn chất và hợp chất.
    }
\end{ex}

\begin{ex}
    Nguyên tố potassium (K) có số hiệu nguyên tử $Z = 19$, thuộc chu kì 4, nhóm IA trong bảng tuần hoàn. Nhận xét nào sau đây về potassium là \textbf{đúng}?
    \choice
    {Potassium là một phi kim hoạt động hóa học mạnh}
    {\True Hydroxide tương ứng với oxide cao nhất của potassium là $\text{KOH}$, có tính base mạnh (dung dịch kiềm)}
    {Oxide cao nhất của potassium có công thức là $\text{KO}_2$ và có tính chất acidic oxide}
    {Potassium có bán kính nguyên tử nhỏ hơn và tính kim loại yếu hơn sodium ($Z = 11$)}
    \loigiai{
        - Do $Z = 19$, K nằm ở chu kì 4, nhóm IA $\rightarrow$ K là một kim loại kiềm điển hình hoạt động rất mạnh.\\
        - Oxide cao nhất là $\text{K}_2\text{O}$ (basic oxide), hydroxide tương ứng là $\text{KOH}$ (kiềm mạnh).\\
        - Trong nhóm IA, đi từ trên xuống dưới theo chiều tăng dần của điện tích hạt nhân, bán kính nguyên tử tăng dần và tính kim loại tăng dần, nên potassium có bán kính lớn hơn và tính kim loại mạnh hơn sodium ($Z = 11$).
    }
\end{ex}

\noindent\textbf{\large PHẦN II. Câu trắc nghiệm đúng sai.} Thí sinh trả lời từ câu 1 đến câu 2. Trong mỗi ý \textbf{a)}, \textbf{b)}, \textbf{c)}, \textbf{d)} ở mỗi câu, thí sinh chọn đúng hoặc sai.

\begin{ex}
    Cho nguyên tử của nguyên tố phosphorus (kí hiệu là P) có số hiệu nguyên tử $Z = 15$.
    \choiceTFt
    {\True Lớp vỏ nguyên tử phosphorus có 3 lớp electron và có 5 phân lớp electron}
    {\True Cấu hình electron của nguyên tử phosphorus được viết dưới dạng thu gọn là $[\text{Ne}]\, 3s^2 3p^3$, nguyên tử có 5 electron ở lớp ngoài cùng}
    {Ở trạng thái cơ bản, nguyên tử phosphorus chỉ có 1 electron độc thân nằm ở phân lớp $3p$}
    {Phosphorus là một nguyên tố phi kim; khi tạo liên kết ion với các kim loại, nguyên tử phosphorus có xu hướng nhường 5 electron để đạt cấu hình bền vững của khí hiếm neon}
    \loigiai{
        \begin{itemchoice}
            \itemch Đúng. Với $Z = 15$, thứ tự mức năng lượng và cấu hình electron là $1s^2 2s^2 2p^6 3s^2 3p^3$. Nguyên tử có 3 lớp electron ($n = 1, 2, 3$) và 5 phân lớp electron gồm $1s, 2s, 2p, 3s, 3p$.
            \itemch Đúng. Viết dưới dạng khí hiếm thu gọn là $[\text{Ne}]\, 3s^2 3p^3$. Lớp ngoài cùng là lớp thứ 3 có $2 + 3 = 5$ electron.
            \itemch Sai. Theo quy tắc Hund, 3 electron trên phân lớp $3p$ phân bố vào 3 orbital riêng biệt ($3p_x^1, 3p_y^1, 3p_z^1$) với spin song song, do đó nguyên tử phosphorus có 3 electron độc thân ở trạng thái cơ bản.
            \itemch Sai. Phosphorus là phi kim có 5 electron lớp ngoài cùng, khi hình thành liên kết ion với kim loại, nguyên tử phosphorus có xu hướng nhận thêm 3 electron để đạt cấu hình octet bền vững của khí hiếm argon gần nhất ($[\text{Ne}]\, 3s^2 3p^6$), tạo thành anion $\text{P}^{3-}$, chứ không nhường 5 electron.
        \end{itemchoice}
    }
\end{ex}

\begin{ex}
    Cho ba nguyên tố $X (Z = 11)$, $Y (Z = 12)$ và $T (Z = 19)$ trong bảng tuần hoàn các nguyên tố hóa học.
    \choiceTFt
    {\True Các nguyên tố $X$ và $Y$ đều thuộc cùng chu kì 3 trong bảng tuần hoàn}
    {\True Cả ba nguyên tố $X, Y, T$ đều thuộc khối các nguyên tố s và đều là các kim loại}
    {\True Giá trị độ âm điện của các nguyên tố giảm dần theo thứ tự: $Y > X > T$}
    {Bán kính nguyên tử của các nguyên tố tăng dần theo thứ tự: $T < X < Y$}
    \loigiai{
        \begin{itemchoice}
            \itemch Đúng. Cấu hình electron của $X (Z = 11)$ là $[\text{Ne}]\, 3s^1$ (chu kì 3, nhóm IA); của $Y (Z = 12)$ là $[\text{Ne}]\, 3s^2$ (chu kì 3, nhóm IIA). Cả hai nguyên tử đều có 3 lớp electron nên thuộc cùng chu kì 3.
            \itemch Đúng. Cấu hình electron của $T (Z = 19)$ là $[\text{Ar}]\, 4s^1$. Electron cuối cùng của cả ba nguyên tử đều điền vào phân lớp s nên $X, Y, T$ đều là các nguyên tố s; các nguyên tử này đều có 1 hoặc 2 electron lớp ngoài cùng nên đều là các kim loại hoạt động mạnh.
            \itemch Đúng. Trong chu kì 3, theo chiều tăng của điện tích hạt nhân thì độ âm điện tăng nên $\chi_Y > \chi_X$ ($\chi_{\text{Mg}} > \chi_{\text{Na}}$). Trong nhóm IA, theo chiều từ trên xuống dưới thì độ âm điện giảm nên $\chi_X > \chi_T$ ($\chi_{\text{Na}} > \chi_{\text{K}}$). Do đó, giá trị độ âm điện giảm dần theo thứ tự: $Y > X > T$.
            \itemch Sai. Trong cùng chu kì 3, từ trái sang phải bán kính nguyên tử giảm: $r_Y < r_X$ ($r_{\text{Mg}} < r_{\text{Na}}$). Trong nhóm IA, từ trên xuống dưới số lớp electron tăng nên bán kính tăng: $r_X < r_T$ ($r_{\text{Na}} < r_{\text{K}}$). Vì vậy, thứ tự tăng dần bán kính nguyên tử đúng phải là $Y < X < T$ ($r_{\text{Mg}} < r_{\text{Na}} < r_{\text{K}}$).
        \end{itemchoice}
    }
\end{ex}

\noindent\textbf{\large PHẦN III. Câu trắc nghiệm trả lời ngắn.} Thí sinh trả lời từ câu 1 đến câu 4.

\begin{ex}
    Cho cấu hình electron lớp ngoài cùng của nguyên tử các nguyên tố sau ở trạng thái cơ bản:
    \begin{itemize}
        \item Nguyên tử A: $3s^2 3p^1$
        \item Nguyên tử B: $3s^2 3p^4$
        \item Nguyên tử C: $2s^2 2p^5$
        \item Nguyên tử D: $4s^1$
        \item Nguyên tử E: $3s^2$
    \end{itemize}
    Trong các nguyên tố trên, có bao nhiêu nguyên tố có tính chất hóa học cơ bản là kim loại?
    \shortans{3}
    \loigiai{
        Dựa vào số electron ở lớp ngoài cùng của nguyên tử để dự đoán tính chất hóa học cơ bản:
        \begin{itemize}
            \item Nguyên tử A có 3 electron lớp ngoài cùng ($3s^2 3p^1$, nguyên tố aluminium) $\rightarrow$ là kim loại.
            \item Nguyên tử B có 6 electron lớp ngoài cùng ($3s^2 3p^4$, nguyên tố sulfur) $\rightarrow$ là phi kim.
            \item Nguyên tử C có 7 electron lớp ngoài cùng ($2s^2 2p^5$, nguyên tố fluorine) $\rightarrow$ là phi kim.
            \item Nguyên tử D có 1 electron lớp ngoài cùng ($4s^1$, nguyên tố potassium) $\rightarrow$ là kim loại.
            \item Nguyên tử E có 2 electron lớp ngoài cùng ($3s^2$, nguyên tố magnesium) $\rightarrow$ là kim loại.
        \end{itemize}
        Vậy có $3$ nguyên tố mang tính chất hóa học cơ bản là kim loại (gồm A, D, E).
    }
\end{ex}

\begin{ex}
    Nguyên tử của nguyên tố $X$ có tổng số hạt cơ bản (proton, neutron, electron) là $34$, trong đó số hạt mang điện nhiều hơn số hạt không mang điện là $10$ hạt. Xác định số thứ tự chu kì của nguyên tố $X$ trong bảng tuần hoàn các nguyên tố hóa học.
    \shortans{3}
    \loigiai{
        Gọi $p, n, e$ lần lượt là số proton, neutron và electron trong nguyên tử $X$. Vì nguyên tử trung hòa về điện nên $p = e = Z$.\\
        Theo bài ra, ta lập hệ phương trình:
        \[
        \begin{cases}
            2Z + n = 34 \\
            2Z - n = 10
        \end{cases}
        \iff
        \begin{cases}
            4Z = 44 \\
            n = 2Z - 10
        \end{cases}
        \iff
        \begin{cases}
            Z = 11 \\
            n = 12
        \end{cases}
        \]
        Với $Z = 11$, cấu hình electron của nguyên tử $X$ là: $1s^2 2s^2 2p^6 3s^1$.\\
        Nguyên tử $X$ có 3 lớp electron, do đó nguyên tố $X$ thuộc chu kì 3 trong bảng tuần hoàn.
    }
\end{ex}

\begin{ex}
    Cho các nguyên tố thuộc chu kì 3 trong bảng tuần hoàn gồm: $_{11}\text{Na}$, $_{12}\text{Mg}$, $_{15}\text{P}$ và $_{17}\text{Cl}$. Trong số các nguyên tố trên, nguyên tố có bán kính nguyên tử lớn nhất có số hiệu nguyên tử $Z$ bằng bao nhiêu?
    \shortans{11}
    \loigiai{
        Các nguyên tố $_{11}\text{Na}$, $_{12}\text{Mg}$, $_{15}\text{P}$ và $_{17}\text{Cl}$ đều thuộc cùng chu kì 3 của bảng tuần hoàn.\\
        Trong một chu kì, đi từ trái sang phải theo chiều tăng dần của điện tích hạt nhân ($Z$ tăng từ $11$ đến $17$), lực hút tĩnh điện giữa hạt nhân và các electron lớp ngoài cùng tăng lên trong khi số lớp electron không đổi, làm cho bán kính nguyên tử giảm dần:
        \[
        r_{\text{Na}} > r_{\text{Mg}} > r_{\text{P}} > r_{\text{Cl}}.
        \]
        Do đó, nguyên tố có bán kính nguyên tử lớn nhất trong dãy là sodium (Na) với số hiệu nguyên tử $Z = 11$.
    }
\end{ex}

\begin{ex}
    Nguyên tố $R$ thuộc nhóm VA trong bảng tuần hoàn các nguyên tố hóa học. Trong hợp chất oxide cao nhất của $R$, nguyên tố oxygen chiếm $56{,}34\%$ về khối lượng. Xác định nguyên tử khối của nguyên tố $R$ (theo đơn vị amu, làm tròn đến hàng đơn vị).
    \shortans{31}
    \loigiai{
        Vì nguyên tố $R$ thuộc nhóm VA nên hóa trị cao nhất của $R$ trong oxide là V.\\
        Công thức oxide cao nhất của $R$ là $\text{R}_2\text{O}_5$.\\
        Thành phần phần trăm khối lượng của oxygen trong hợp chất oxide cao nhất là:
        \[
        \%m_{\text{O}} = \dfrac{5 \times 16}{2 \times M_R + 5 \times 16} \times 100\% = 56{,}34\%
        \]
        \[
        \iff \dfrac{80}{2M_R + 80} = 0{,}5634 \implies 2M_R + 80 = \dfrac{80}{0{,}5634} \approx 141{,}995
        \]
        \[
        \implies 2M_R \approx 61{,}995 \implies M_R \approx 30{,}997 \approx 31\text{ amu}.
        \]
        Vậy nguyên tử khối của nguyên tố $R$ là $31$ (nguyên tố phosphorus, P).
    }
\end{ex}
